%% OVERFORE-Formulario-T-18.tex — generado desde OVERFORE-Formulario-T-18.docx
%% Formato: clase overfore, modo formulario. El formato lo pone la clase; aquí sólo va el contenido.
\documentclass[formulario]{../../comun/overfore}
\numsubdoc{7}
\formulario{T-18}{Propuesta de implantación y puesta en marcha controlada}
\informe{Informe Preparatorio N° 2}
\fecha{12 de octubre de 2026}
\externos{T14}
\begin{document}

\apartado{Estrategia de puesta en marcha}
Este formulario presenta la implantación y la puesta en marcha controlada de SITCO Aranda que piden las Bases administrativas para el Formulario T-18 (Compañía Minera Altos de Aranda [CMAA], 2026a, p. 64): la Etapa 1, con su marcha blanca de los meses 13 a 15 y su producción desde el mes 16, y la Etapa 2, con su marcha blanca de los meses 19 y 20 y su producción desde el mes 21, más la convivencia entre ambas y el procedimiento de reversión. Las fechas, dotaciones y criterios son los de las fichas del Formulario T-14 (Overfore S.A., 2026a), y sus horas, las del Formulario T-15 (Overfore S.A., 2026b).

Cada etapa entra a producción por área, en tres olas, después de una marcha blanca en la que SITCO convive con la forma actual de trabajar y se concilia con ella cada día. El orden de las olas es el del Subdocumento 3, sección 3.2.2.5 (Overfore S.A., 2026c): en la Etapa 1, portería y acreditación, mina y al final planta, concentrado y puerto, en la Etapa 2, seguridad y salud ocupacional, equipos y mantenimiento y al final cumplimiento ambiental con el depósito de relaves. La Figura \ref{fig:calendario-implantacion-dos-etapas-meses-12} muestra el calendario, con el mes 1 en enero de 2027 (CMAA, 2026a, Art. 17.1, p. 12).

\figura{16.0}{fig-A-1-calendario-implantacion-dos-etapas-meses}{Calendario de implantación de las dos etapas, meses 12 a 23.}{fig:calendario-implantacion-dos-etapas-meses-12}{Fuente: elaboración propia con las fichas del Formulario T-14.}
La figura muestra que ninguna ola cae en una detención de planta: las de marzo y septiembre de 2028 se usan para el corte de 48 horas y la conmutación de recuperación de cada etapa, antes de su paso. Cada área se capacita durante dos ciclos del 7x7 antes de entrar a su marcha blanca, y entre el cierre de la marcha blanca y la ola de cada área la convivencia continúa, porque las planillas dejan de ser fuente oficial sólo el día del paso. La siguiente tabla muestra cómo el plan cumple las seis condiciones que el cliente impone a toda estrategia de puesta en producción (CMAA, 2026b, 13.3, p. 22).

\begin{tabla}{I{0.330}I{0.451}I{0.220}}{Condiciones del cliente para la puesta en producción y cómo se cumplen.}{tab:condiciones-cliente-puesta-produccion-como-se}
  {\cab{Condición del cliente} & \cab{Cómo la cumple el plan} & \cab{Paquetes del T-14}}
  {Fuente: elaboración propia a partir de las Bases técnicas del caso Minería (CMAA, 2026b, cap. 17, pp. 29–31) y de las Bases administrativas (CMAA, 2026a, cap. 3, Art. 17, pp. 12–13).}
1. Convivir con la forma actual, conciliar y poder volver atrás & Marcha blanca con conciliación diaria contra las planillas y reversión por área en 30 minutos o menos & 1.6.3.1, 1.6.3.2, 1.6.4.1 y 1.6.4.2 \\
2. Ningún paso en una detención ni en el cierre contable & Cada ola va después del quinto día hábil, en la ventana de 03:00 a 05:00 y fuera de las 48 horas previas a un embarque y del embarque mismo & 1.6.4.1 y 1.6.4.2 \\
3. Despliegue por área o sitio & Tres olas por etapa, separadas por al menos 72 horas & 1.6.4.1 y 1.6.4.2 \\
4. Acompañamiento de día y de noche en cada ola & Una persona por turno en cada área en marcha blanca y en sus dos primeras semanas de producción & 1.6.3.3, 1.6.3.4, 1.6.4.1 y 1.6.4.2 \\
5. Dos ciclos del 7x7 para capacitar un área & Cuatro semanas de capacitación por ola, en los dos grupos y en ambos turnos & 1.6.5.1 y 1.6.6.5 \\
6. Estabilización con dotación y duración declaradas, en faena & Hasta el 31 de mayo de 2028 en la Etapa 1 y hasta el 8 de noviembre de 2028 en la Etapa 2, sin costo adicional & 1.6.4.1, 1.6.4.2 y 1.6.4.3 \\
\end{tabla}

\apartado{Implantación por sitio y perfil}
El plan de implantación se organiza por sitio y por perfil, como exige el RT-20.01 (CMAA, 2026c, p. 35). Cada ola corresponde a un sitio, con sus canales y sus equipos, y se capacitan los cinco perfiles del plan de capacitación: personas usuarias finales, usuarias avanzadas, administradoras, equipo técnico y soporte de niveles 1 y 2. La siguiente tabla los resume.

\begin{tabla}{I{0.253}I{0.275}I{0.242}I{0.231}}{Sitios, olas, canales y perfiles de la implantación.}{tab:sitios-olas-canales-perfiles-implantacion}
  {\cab{Sitio y ola} & \cab{Canales y equipos} & \cab{Perfiles principales} & \cab{Población}}
  {Fuente: elaboración propia a partir del Subdocumento 3 (sección 3.3) y de las Bases técnicas del caso Minería (CMAA, 2026b, cap. 3, p. 6).}
Portería y acreditación (Etapa 1, ola 1) & Consola de Portería en dos PC industriales y Portal de Contratistas & Personal de portería y de acreditación, empresas contratistas & 4.180 personas con acceso, hasta 8.500 en detención, 2.900 pasos diarios \\
Mina (Etapa 1, ola 2) & Botonera de Pala en las cinco palas, con una pala piloto de 14 días, y 10 tabletas & Operadores de pala, supervisores de turno y despachadores & 20 operadores de pala (5 palas × 2 turnos × 2 grupos del 7x7) y 10 puestos de supervisión \\
Planta, concentrado y puerto (Etapa 1, ola 3) & 6 tabletas en planta y 2 en el acopio de puerto, con su gabinete de borde & Supervisores de planta, personal del acopio y del embarque & 8 puestos, 12.400 viajes de concentrado y 31 embarques al año \\
Seguridad y salud Ocupacional (Etapa 2, ola 1) & 4 tabletas, lectores faciales en cinco carriles de portería, sólo con acuerdo sindical & Inspectores y supervisores de seguridad & 4 puestos, 7.800 inspecciones al año \\
Equipos y mantenimiento (Etapa 2, ola 2) & 8 tabletas & Mantenedores y planificadores & 8 puestos, 34.000 órdenes de trabajo al año \\
Cumplimiento ambiental (Etapa 2, ola 3) & Caudalímetros de planta, campamento y carga de aljibes, y gabinete de borde del depósito de relaves & Encargados ambientales & 3 caudalímetros y 186 instrumentos \\
Área de TI del cliente (ambas etapas) & Consolas de administración y siete estaciones & Administradores y equipo técnico & 11 personas, 4 de ellas en faena \\
\end{tabla}
La tabla muestra que ninguna ola reúne más de un sitio de la faena, salvo la tercera de la Etapa 1, que suma el acopio de puerto porque cierra la cadena del lote embarcado y tiene su propia persona en Antofagasta. La población más grande es la de portería, y entra primero porque usa la credencial que ya tienen todas las personas con acceso. Los equipos de terreno los compra el cliente con la especificación de Overfore y quedan instalados antes de la marcha blanca de su etapa: los de la Etapa 1 antes del mes 13 y los de la Etapa 2 antes del mes 19. Las fichas 1.6.2.1 y 1.6.2.2 del Formulario T-14 fijan la compra: seis botoneras, 18 tabletas y dos equipos de portería en la Etapa 1, y siete lectores faciales, cuatro caudalímetros y 15 tabletas en la Etapa 2. Lo que excede los equipos en uso de la tabla queda como repuesto.

\apartado{Etapa 1: marcha blanca de los meses 13 a 15 y producción desde el mes 16}
La marcha blanca de la Etapa 1 dura los tres meses del Art. 17.1, del 1 de enero al 31 de marzo de 2028, con datos y usuarios reales desde el primer día. Las áreas entran escalonadas, cada una después de sus cuatro semanas de capacitación, y pasan a producción en abril, en las ventanas que siguen al quinto día hábil, el viernes 7. La siguiente tabla muestra las tres olas.

\begin{tabla}{I{0.209}I{0.187}I{0.253}I{0.176}I{0.176}}{Olas de la Etapa 1.}{tab:olas-etapa-1}
  {\cab{Ola} & \cab{Capacitación} & \cab{Entrada a marcha blanca} & \cab{Paso a producción} & \cab{Estabilización por área}}
  {Fuente: elaboración propia a partir de las Bases administrativas (CMAA, 2026a, cap. 3, Art. 17, pp. 12–13); Subdocumento 7, cap. 7, sección 7.3, p. 9.}
1 · Portería y acreditación & Diciembre de 2027 & 1 de enero de 2028 & Lunes 10 de abril & Hasta el 23 de abril \\
2 · Mina & Diciembre de 2027 y enero de 2028 & 15 de enero, con la pala piloto, las demás palas desde la quinta semana & Lunes 17 de abril & Hasta el 30 de abril \\
3 · Planta, concentrado y puerto & Enero de 2028 & 5 de febrero & Viernes 21 de abril & Hasta el 4 de mayo \\
\end{tabla}
Una ola pasa sólo si la anterior cumplió 72 horas sin incidentes críticos ni altos, no tiene diferencias no explicadas en su conciliación y alcanzó su meta de adopción: 95 \% o más de los registros de su área hechos en SITCO y, en mina, 90 \% de las declaraciones hechas con la Botonera por el propio operador. Si no la alcanza, la ola siguiente pasa a la siguiente ventana que cumpla las reglas, y quedan ventanas hasta el 28 de abril sin mover H7. La Botonera se extiende a las demás palas sólo con el acta de la gerencia de operaciones mina, después de catorce días de piloto con los dos grupos del 7x7 y ambos turnos.

Tres actividades fijan la entrada a producción. Los datos de la Etapa 1 se cargan al final del mes 12, la carga incremental trae lo que las planillas registran durante la marcha blanca y el acta de conciliación se firma al cierre del mes 15. En la detención de planta de marzo de 2028, del 13 al 18, se ejecutan el corte controlado de 48 horas y la conmutación a la región de recuperación. La conmutación verifica un RTO de 4 horas y un RPO de 15 minutos, y el corte, cero registros perdidos y sincronización en 4 horas o menos. Cada ola la aprueba el Comité de Proyecto, el Comité de Operación autoriza su intervención menor, y el acta de paso, firmada en el mes 16, es la evidencia de H7.

\apartado{Indicadores diarios y cierre de cada marcha blanca}
Una marcha blanca se cierra sólo cuando se cumplen las seis condiciones del Art. 17.3 (CMAA, 2026a, p. 13), medidas cada día durante al menos sus cuatro últimas semanas, como pide el RT-20.04 (CMAA, 2026c, p. 35). La siguiente tabla muestra el indicador diario y el umbral de cierre de cada condición.

\begin{tabla}{I{0.297}I{0.319}I{0.385}}{Indicadores diarios y umbrales de cierre de las marchas blancas.}{tab:indicadores-diarios-umbrales-cierre-marchas-blancas}
  {\cab{Condición del Art. 17.3} & \cab{Indicador diario} & \cab{Umbral de cierre}}
  {Fuente: elaboración propia a partir de las Bases administrativas (CMAA, 2026a, cap. 3, Art. 17, pp. 12–13); Subdocumento 9, cap. 9, sección 9.2.6, p. 21.}
Sin incidentes críticos ni altos & Incidentes abiertos atribuibles a la solución, por severidad & Cero críticos y cero altos. Defectos medios y bajos: hasta 4 y 13 en la Etapa 1, y hasta 1 y 5 en la Etapa 2, con plan y fecha \\
Volumen real comprometido & Proporción de los registros del período hechos en SITCO & 95 \% o más: viajes, pesajes, lotes y pasos de portería en la Etapa 1, avisos de falla, registros de seguridad y mediciones ambientales en la Etapa 2 \\
Disponibilidad y tiempo de respuesta & Disponibilidad de los servicios críticos y tiempo del registro en terreno & 99,9 \% o más y 2 segundos o menos en el percentil 95 \\
Conciliación sin diferencias no explicadas & Diferencias entre SITCO y las planillas o registros vigentes & Ninguna diferencia sin explicar \\
Personal capacitado y certificado & Personas de las olas capacitadas, por área y turno & 100 \% conforme al plan de capacitación aprobado \\
Acta de aceptación & Avance del expediente del hito & Acta suscrita por la Contraparte Técnica \\
\end{tabla}
Además, en la Etapa 1 se mide en cada turno el tiempo que agrega al operador de pala cada declaración y cada punto de muestreo, que no puede superar el que justifica el diseño, porque la solución no puede agregar segundos al ciclo. Como las cuatro últimas semanas de la Etapa 1 incluyen la detención de marzo de 2028, el volumen comprometido de la ola de planta, concentrado y puerto descuenta sus seis días, y en ellos la conciliación sigue con las planillas, sin validaciones con el personal de planta. Si al cierre una condición no se cumple, Overfore extiende la marcha blanca a su costo, sin mover las fechas de las fases siguientes y afecto a las multas del Art. 80.

\apartado{Etapa 2: marcha blanca de los meses 19 y 20 y producción desde el mes 21}
La marcha blanca de la Etapa 2 ocupa julio y agosto de 2028. Sus áreas entran en semanas sucesivas, cada una después de sus cuatro semanas de capacitación, y pasan a producción después de la detención de planta de septiembre, del 11 al 16, en la que se repiten el corte de 48 horas y la conmutación de recuperación. La siguiente tabla muestra las tres olas.

\begin{tabla}{I{0.253}I{0.176}I{0.198}I{0.198}I{0.176}}{Olas de la Etapa 2.}{tab:olas-etapa-2}
  {\cab{Ola} & \cab{Capacitación} & \cab{Entrada a marcha blanca} & \cab{Paso a producción} & \cab{Estabilización por área}}
  {Fuente: elaboración propia a partir de las Bases administrativas (CMAA, 2026a, cap. 3, Art. 17, pp. 12–13); Subdocumento 7, cap. 7, sección 7.3, p. 9.}
1 · Seguridad y salud ocupacional & Junio de 2028 & Primera semana de julio & Miércoles 20 de septiembre & Hasta el 3 de octubre \\
2 · Equipos y mantenimiento & Junio y julio de 2028 & Segunda semana de julio & Domingo 24 de septiembre & Hasta el 7 de octubre \\
3 · Cumplimiento ambiental & Junio y julio de 2028 & Tercera semana de julio & Jueves 28 de septiembre & Hasta el 11 de octubre \\
\end{tabla}
Las olas pasan después del quinto día hábil, el jueves 7, y de las fiestas del 18 y 19, y quedan el 29 y el 30 como ventanas de reserva para H12. Los datos de la Etapa 2 se cargan al final del mes 18 y su acta de conciliación se firma al cierre del mes 20, que incluye los formularios de seguridad digitalizados. El expediente de H12 se entrega al cierre del mes 20, para que los diez días hábiles de revisión corran antes de las olas. Sin acuerdo con el sindicato al inicio del mes 15, el cliente no compra los lectores faciales y la portería sigue con la credencial de proximidad. Con acuerdo, la verificación facial entra con la primera ola, después de la detención de septiembre, comenzando por el personal permanente y con la credencial como alternativa, el personal de detención se incorpora desde la detención de marzo siguiente, con enrolamiento anticipado por empresa contratista. Los reportes ambientales se generan en paralelo a los manuales durante la marcha blanca.

\apartado{Convivencia entre las dos etapas}
Entre los meses 19 y 20 la Etapa 2 en marcha blanca convive con la Etapa 1 en producción, y su paso no puede degradar a la Etapa 1 (CMAA, 2026a, Art. 17.2, p. 13). La siguiente tabla muestra qué es fuente oficial en cada período, cómo se concilia y cuándo se apaga lo anterior.

\begin{tabla}{I{0.297}I{0.231}I{0.231}I{0.242}}{Fuente oficial, conciliación y apagado de lo anterior por período.}{tab:fuente-oficial-conciliacion-apagado-lo-anterior}
  {\cab{Período} & \cab{Fuente oficial} & \cab{Conciliación diaria} & \cab{Apagado de lo anterior}}
  {Fuente: elaboración propia a partir de las Bases administrativas (CMAA, 2026a, cap. 3, Art. 17, pp. 12–13); Subdocumento 7, cap. 7, sección 7.3, p. 9.}
Marcha blanca de la Etapa 1 (enero a marzo de 2028) & Planillas vigentes de cada área & SITCO contra las planillas & El día del paso de cada área, sus planillas quedan archivadas en sólo lectura \\
Etapa 1 en producción y desarrollo de la Etapa 2 (abril a junio de 2028) & SITCO para la Etapa 1 & Sin conciliación con planillas & Ya apagado \\
Marcha blanca de la Etapa 2 (julio y agosto de 2028, pasos de cada área en septiembre) & SITCO para la Etapa 1, planillas sólo para los procesos de la Etapa 2 & SITCO contra los registros de cada área de la Etapa 2 & El día del paso de cada área de la Etapa 2 \\
Operación (desde el 20 de septiembre de 2028) & SITCO para ambas etapas & Sin conciliación con planillas & Planillas archivadas \\
\end{tabla}
Los datos que comparten ambas etapas tienen una sola fuente y no se digitan dos veces. Para no degradar la Etapa 1, los módulos de la Etapa 2 se despliegan como servicios separados y cada promoción ejecuta la regresión de la Etapa 1. Entre junio y agosto de 2028 una persona por turno atiende la Etapa 1 en producción, sin incidentes críticos ni altos abiertos al inicio de la operación, y los puestos de atención continua de la operación se dotan desde el 1 de septiembre.

\apartado{Procedimiento de reversión}
El paso a producción es gradual y reversible, como exige el RT-20.02 (CMAA, 2026c, p. 35): se revierte una función o un área, sin detener las demás. La siguiente tabla muestra el procedimiento, quién lo decide y cuánto demora.

\begin{tabla}{I{0.187}I{0.407}I{0.220}I{0.187}}{Procedimiento de reversión de una ola.}{tab:procedimiento-reversion-ola}
  {\cab{Paso} & \cab{Qué se hace} & \cab{Responsable} & \cab{Tiempo}}
  {Fuente: elaboración propia a partir de Subdocumento 7, cap. 7, sección 7.3, p. 9.}
Disparador & Un incidente crítico no resuelto en su plazo de 4 horas durante la ventana o las 72 horas siguientes, o diferencias no explicadas en la conciliación diaria & Líder de Operación / SRE informa & Al detectarse \\
Decisión & El Comité de Operación ordena revertir la función o el área afectada & Comité de Operación & En la misma ventana o turno \\
Ejecución & Se reactiva la versión anterior con el flujo de entrega, y las planillas y la radio del área vuelven a ser la forma oficial de trabajo & Líder de Operación / SRE & 30 minutos o menos desde la decisión, dentro de la ventana de dos horas \\
Recuperación de datos & Las planillas reciben los registros que SITCO tomó en el período antes de reanudar & Líder de Implantación y Gestión del Cambio & Antes de reanudar la operación del área \\
Nuevo intento & La ola vuelve a pasar en la siguiente ventana que cumpla 72 horas sin incidentes críticos ni altos & Comité de Proyecto & Siguiente ventana válida \\
\end{tabla}
Al revertir no se pierden registros, porque el Nodo de Faena conserva su bandeja de salida y lo registrado en SITCO pasa a las planillas. Lo que se pierde es la ventana: el área vuelve a operar con planillas y radio hasta su nuevo paso. La reversión se ensaya en Preproducción con su tiempo medido antes de cada paso, y se ensaya otra vez durante cada marcha blanca, un plan de reversión activo es condición para iniciar la marcha blanca de la Etapa 1 (H6).

\apartado{Acompañamiento, estabilización y adopción}
El acompañamiento en terreno tiene dotación declarada y decreciente según la adopción, y cada paso tiene una estabilización con atención reforzada y sin costo adicional, como piden los RT-20.05 y RT-20.06 (CMAA, 2026c, p. 35). La siguiente tabla resume la dotación, en turnos de doce horas de día y de noche.

\begin{tabla}{I{0.319}I{0.319}D{0.132}D{0.110}I{0.121}}{Dotación de acompañamiento y estabilización en faena.}{tab:dotacion-acompanamiento-estabilizacion-faena}
  {\cab{Período} & \cab{Personas por turno} & \cab{Persona-turnos} & \cab{HH} & \cab{Paquete}}
  {Fuente: elaboración propia a partir de Formulario T-14 (asociado al cap. 7), tabla \refext{T14}{tab:esfuerzo-estimado-rama-origen-horas} del Formulario T-14 (esfuerzo estimado por rama), p. 45.}
Marcha blanca de la Etapa 1, 13 semanas & Una en las semanas 1 y 2, dos en las 3 a 5 y cuatro en las 6 a 13 & 560 & 6.824 & 1.6.3.3 \\
Estabilización de la Etapa 1, del 10 de abril al 31 de mayo de 2028 & Una por área en sus dos primeras semanas, después, una para toda la faena & 166 & 2.064 & 1.6.4.1 \\
Soporte de la Etapa 1, junio a agosto de 2028 & Una, con especialistas de nivel 2 en ambos turnos & 184 & 2.280 & 1.6.4.3 \\
Marcha blanca de la Etapa 2, 9 semanas & Una en la semana 1, dos en la 2 y tres desde la 3 & 330 & 4.248 & 1.6.3.4 \\
Estabilización de la Etapa 2, del 20 de septiembre al 8 de noviembre de 2028 & Una por área en sus dos primeras semanas, después, una para toda la faena & 140 & 1.752 & 1.6.4.2 \\
\end{tabla}
Una persona por turno basta en las semanas 1 y 2 de la marcha blanca de la Etapa 1 porque en ellas sólo opera portería, que entra primero y usa la credencial de proximidad que ya tienen todas las personas con acceso: los unos 2.900 pasos diarios no piden un gesto nuevo, sino acercar la misma tarjeta a un lector distinto. La persona de cada turno se ubica en portería en los cambios de turno de las 08:00 y las 20:00, cuando se concentra el flujo, y la mesa de ayuda 24x7 atiende el resto del día. Si en un cambio de turno la cola supera diez personas o el tiempo de paso supera el umbral de RNF-DES-003, se suma un técnico de terreno acreditado al turno siguiente con horas de la contingencia de implantación del Subdocumento 8, sin cambiar la dotación planificada.

La adopción se mide cada mes por área y grupo del 7x7, con fuentes independientes del registro, como la telemetría, el despacho, la báscula y la portería. La meta al cierre de cada marcha blanca es 90 \% de los cambios de material declarados por el propio operador con la Botonera, con guantes y sin registro paralelo, y ninguna planilla paralela en las áreas que ya pasaron a producción. Un área bajo la meta recibe en dos semanas un plan de acción con acompañamiento adicional en su turno. El dato por persona sólo orienta el acompañamiento y no se usa para evaluar ni sancionar a nadie. Quienes acompañan trabajan en el régimen 7x7, se alojan en el campamento a 9 km de la faena y, como cualquier contratista, tienen exámenes preocupacionales y de altura, inducción de seguridad y credencial vigente antes de su primer turno.

Después del paso de la Etapa 2, los meses 22 a 27 se dedican a la mentoría de las once personas del área de TI del cliente, con cada competencia certificada con un ejercicio práctico. La mesa de ayuda, los centros de operaciones de red y de seguridad, las pruebas periódicas y la recuperación ante desastres quedan como servicio de Overfore durante la operación. En el mes 21 la operación llega a 38 personas, por la estabilización de la Etapa 2, y en régimen tiene entre 30 y 31 (Overfore S.A., 2026b). De ellas, 20 cubren los cuatro puestos de atención continua, con cinco personas por puesto para las 24 horas en el 7x7.

\seccionsinnumero{Referencias}
\begin{referencias}
\item Compañía Minera Altos de Aranda. (2026a). \emph{Bases administrativas}. Licitación N° TFEP-01/2026.
\item Compañía Minera Altos de Aranda. (2026b). \emph{Bases técnicas del caso Minería}. Licitación N° TFEP-01/2026.
\item Compañía Minera Altos de Aranda. (2026c). \emph{Bases técnicas transversales}. Licitación N° TFEP-01/2026.
\item Overfore S.A. (2026a). Formulario T-14: Plan de trabajo, EDT y carta Gantt.
\item Overfore S.A. (2026b). Formulario T-15: Nivelación de recursos.
\item Overfore S.A. (2026c). Subdocumento 3: Esquema de solución y alcance. Licitación TFEP-01/2026.
\end{referencias}
\end{document}
